\ssr{ВВЕДЕНИЕ}
%<цель и задачи; можно вводное слово к работе>

Цель работы -- разработать и реализовать программное обеспечение на языке python для извлечения данных из текстовых файлов, полученных из pdf-файлов с применением библиотеки PyPDF2, с использованием регулярных выражений.

Задачи, необходимые для достижения поставленной цели:
\begin{enumerate}
	\item разработать регулярные выражения для решения задачи по варианту;
	\item реализовать функцию для поиска подстроки по варианту в pdf-файле с использованием разработанных регулярных выражений (привести листинг как самой функции, так и списка необходимых для её функционирования библиотек, а также пример вызова функции), возвращающую кортеж вида: 
	\begin{itemize}
		\item значение, определяющее наличие строки согласно заданному регулярному выражению ("истина/ложь"); 
		\item список кортежей (если "ложь", то список пустой) с двумя составляющими каждого кортежа --- строкой, найденной с помощью регулярного выражения, и координатами для определения найденной строки в документе (минимально необходимо указать номер страницы и номер строки документа, в которой найдено вхождение искомой строки. непосредственно таблицы и изображения строками не считать, но их названия считаются строками);
	\end{itemize}
	\item реализовать программное обеспечение, принимающее на вход путь к pdf-файлу, использующее функцию из предыдущего пункта;
	\item проверить реализацию на приложенных к лабораторной работе файлах, привести таблицу с колонками "название использованного pdf-файла", "признак успешного нахождения подстроки", "координаты первого нахождения подстроки". необходимо приложить файлы, в которых были найдены ошибки, в репозиторий лабораторной работы.
\end{enumerate}

% Требуется достичь поставленной цели с использованием не менее трёх больших языковых моделей (желательно локально разворачиваемых, например, qwen3-coder, https://ollama.com/search?q=coder или аналогичных). В отчёте отобразить все попытки.\\
% Очевидно:
% \begin{itemize}
% 	\item в аналитическом разделе указать краткое описание используемых БЯМов;
% 	\item в конструкторском разделе указать описание процесса взаимодействия с каждой выбранной БЯМ в рамках решения задачи;
% 	\item в технологическом разделе указать ПРОМПТы, их результаты, внесённые правки, результаты правок и т. д. до получения итогового результата;
% 	\item в исследовательском разделе отобразить полученные результаты для приложенного набора файлов, оценить сложность полученного силами БЯМ алгоритма решения задачи.
% \end{itemize}

\clearpage
